% Part:sets-functions-relations
% Chapter: size-of-sets
% Section: reduction-alt

\documentclass[../../../include/open-logic-section]{subfiles}

\begin{document}

\olfileid{sfr}{siz}{red-alt}

\olsection{Reduksi}

\begin{editorial}
  Bagean iki mbuktekake sipat ora bisa didhaptar kanthi reduksi,
  cocog karo asil ing \olref[nen-alt]{sec}. Versi alternatif
  kang rada luwih rinci, cocog karo asil ing \olref[nen]{sec},
  diwenehake ing \olref[red]{sec}.
\end{editorial}

Kita wis mbuktekake yen $\Bin^\omega$ iku !!{nonenumerable} nganggo
argumen diagonalisasi. Argumen diagonalisasi kang padha digunakake
kanggo nuduhake yen $\Pow{\Nat}$ iku !!{nonenumerable}.
Nanging ana cara liya kanggo mbuktekake yen $\Pow{\Nat}$ iku !!{nonenumerable}:
tuduhna yen \emph{yen $\Pow{\Nat}$ iku !!{enumerable}, mula $\Bin^\omega$
uga !!{enumerable}}. Amarga kita ngerti yen $\Bin^\omega$
iku !!{nonenumerable}, mula $\Pow{\Nat}$ uga mangkono.

Iki diarani \emph{ngreduksi} sawijining masalah menyang masalah liyane.
Ing kene, masalah ngenumerasi $\Bin^\omega$ kita reduksi menyang
masalah ngenumerasi $\Pow{\Nat}$. Solusi kanggo masalah kang pungkasan,
yaiku enumerasi $\Pow{\Nat}$, bakal ngasilake solusi
kanggo masalah kapisan, yaiku enumerasi $\Bin^\omega$.

Kanggo ngreduksi masalah ngenumerasi himpunan~$B$ menyang masalah ngenumerasi
himpunan~$A$, kita menehi cara ngowahi enumerasi saka~$A$
dadi enumerasi saka~$B$. Cara kang paling gampang yaiku netepake
!!a{surjection} $f\colon A \to B$. Yen $x_1$, $x_2$, \dots{}
ngenumerasi~$A$, mula $f(x_1)$, $f(x_2)$, \dots{} ngenumerasi~$B$.
Ing kasus iki, kita nggoleki !!a{surjection} $f\colon \Pow{\Nat}
\to \Bin^\omega$.

\begin{prob}
Tuduhna yen ana fungsi !!{injective} $g\colon B \to A$, lan
$B$~iku !!{nonenumerable}, mula~$A$ uga mangkono. Tuduhna carane
nggunakake~$g$ kanggo ngowahi enumerasi saka~$A$ dadi enumerasi saka~$B$.
\end{prob}

\begin{proof}[Bukti {\olref[nen-alt]{thm:nonenum-pownat}} kanthi reduksi]
Kanggo reduksi, upamakna $\Pow{\Nat}$ iku !!{enumerable}, mula
ana enumerasi $N_{1}$, $N_{2}$, $N_{3}$, \dots

Tetepna fungsi $f \colon \Pow{\Nat} \to \Bin^\omega$ kanthi netepake
$f(N)$ dadi rerangken $s_{k}$ saengga $s_{k}(n) = 1$ yen lan mung yen $n \in N$,
lan $s_k(n) = 0$ yen ora mangkono.

Iki cetha netepake sawijining fungsi, amarga yen $N \subseteq \Nat$,
saben $n \in \Nat$ mesthi !!a{element} saka $N$ utawa dudu anggotane.
Upamane, himpunan $2\Nat = \Setabs{2n}{n \in \Nat} = \{0,2, 4, 6,
\dots\}$ saka wilangan natural genep dipetakake menyang rerangken $1010101\dots$;
$\emptyset$ dipetakake menyang $0000\dots$; $\Nat$ dipetakake menyang
$1111\dots$.

Fungsi iki uga !!{surjective}: saben rerangken $0$ lan $1$ cocog
karo sawijining himpunan wilangan natural, yaiku himpunan kang
anggotane wilangan natural kang nuduhake posisi rerangken kang isine~$1$.
Kanthi luwih pas, yen $s \in \Bin^\omega$, tetepna $N
\subseteq \Nat$ kanthi:
\[
N = \Setabs{n \in \Nat}{s(n) = 1}
\]
Dadi $f(N) = s$, kaya kang bisa dipriksa saka definisi~$f$.

Saiki gatekna dhaptar
\[
f(N_1), f(N_2), f(N_3), \dots
\]
Amarga $f$ iku !!{surjective}, saben anggota $\Bin^\omega$ mesthi
katon minangka nilai~$f$ kanggo sawijining argumen, mula mesthi
katon ing dhaptar. Dadi dhaptar iki ngenumerasi kabeh~$\Bin^\omega$.

Mula yen $\Pow{\Nat}$ iku !!{enumerable}, $\Bin^\omega$ mesthi
!!{enumerable}. Nanging $\Bin^\omega$ iku !!{nonenumerable}
(\olref[nen-alt]{thm:nonenum-bin-omega}). Dadi $\Pow{\Nat}$ iku
!!{nonenumerable}.
\end{proof}

%\begin{explain}
%Arah reduksi gampang mbingungake.
%Upamane, fungsi !!a{surjective} $g \colon \Bin^\omega \to X$
%\emph{ora} mbuktekake yen $X$ iku !!{nonenumerable}. (Gatekna $g
%\colon \Bin^\omega \to \Bin$ kang ditetepake dening $g(s) = s(1)$,
%yaiku fungsi kang metakake rerangken $0$ lan $1$ menyang !!{element} kapisane.
%Fungsi iki surjektif, amarga ana rerangken kang diwiwiti $0$ lan
%ana kang diwiwiti $1$. Nanging $\Bin$ winates.) Gatekna uga yen fungsi $f$
%kudu surjektif; yen ora, argumen iki ora lumaku:
%$f(x_1)$, $f(x_2)$, \dots{} ora mesthi ngemot
%kabeh !!{element}s saka~$Y$. Upamane, $h\colon \Nat \to
%\Bin^\omega$ kang ditetepake dening
%\[
%h(n) = \underbrace{000\dots0}_{\text{$n$ wilangan $0$}}
%\]
%iku fungsi, nanging $\Nat$ iku !!{enumerable}.
%\end{explain}

\begin{prob}\label{sfr:siz:red:prob:nat-nat}
  Tuduhna yen himpunan~$X$ saka kabeh fungsi $f\colon \Nat \to \Nat$
  iku !!{nonenumerable} nganggo argumen reduksi. (Pituduh: wenehana fungsi
  surjektif saka $X$ menyang~$\Bin^\omega$.)
\end{prob}

\begin{prob}
Tuduhna yen himpunan kabeh \emph{himpunan saka} pasangan wilangan natural,
yaiku $\Pow{\Nat \times \Nat}$, iku !!{nonenumerable} nganggo argumen
reduksi.
\end{prob}

\begin{prob}
Tuduhna yen $\Nat^\omega$, himpunan rerangken tanpa wates saka wilangan
natural, iku !!{nonenumerable} nganggo argumen reduksi.
\end{prob}

%\begin{prob}
%Upamakna $P$ himpunan fungsi saka $\Nat$ menyang himpunan $\{0\}$, lan $Q$ himpunan fungsi \emph{parsial}
%saka himpunan wilangan wutuh positif menyang himpunan $\{0\}$. Tuduhna
%yen $P$~iku !!{enumerable} lan $Q$~ora. (Pituduh: reduksinen masalah
%ngenumerasi $\Bin^\omega$ menyang ngenumerasi~$Q$.)
%\end{prob}

\begin{prob}
Upamakna $S$ himpunan kabeh !!{surjection}s saka $\Nat$ menyang himpunan
$\{0,1\}$, yaiku $S$ ngemot kabeh !!{surjection}s~$f \colon \Nat
\to \Bin$. Tuduhna yen $S$ iku !!{nonenumerable}.
\end{prob}

\begin{prob}
Tuduhna yen himpunan~$\Real$ saka kabeh wilangan real iku !!{nonenumerable}.
\end{prob}

\end{document}
